% !TEX root = ../master.tex

\section*{Syllabus}

This course provides an introduction to the physical and mechanical processes governing earthquake generation and seismic radiation. Starting from continuum mechanics and elastodynamics, students develop the mathematical tools needed to describe stress, strain, elastic waves, Green's functions, and seismic sources. The course then examines moment tensors, radiation patterns, seismic moment, source spectra, and kinematic source models. Building on these foundations, it introduces fracture mechanics, friction laws, laboratory observations, and the dynamics of earthquake rupture, including supershear and slow-slip phenomena. Particular attention is given to the energy budget of earthquakes and to the connection between individual earthquake sources and statistical patterns of seismicity. The course combines theoretical developments, classical models, observational evidence, and laboratory experiments to provide a unified framework for understanding the earthquake source process.

%\bigskip

%\bigskip
%
%\begin{longtable}{|p{0.5cm}|c|p{6cm}|c|}
%\hline
% & \textbf{Date} & \textbf{Topic} & \textbf{Instructor} \\
%\hline
%\endhead
%1  & Fri 11 Sep 2026 & Continuum Mechanics \& Math Primer & Harsha \\
%\hline
%2  & Fri 18 Sep 2026 & Elastodynamics I: The Wave Equation & Harsha \\
%\hline
%3  & Fri 09 Oct 2026 & Elastodynamics II: Green's Functions \& Representation Theorem & Harsha \\
%\hline
%4  & Fri 16 Oct 2026 & Far-Field Radiation \& the Seismic Source & Pierre \\
%\hline
%5*  & Fri 23 Oct 2026 & Kinematic Source Models & Pierre \\
%\hline
%6*  & Fri 30 Oct 2026 & Fracture Mechanics Fundamentals & Harsha \\
%\hline
%7  & Fri 06 Nov 2026 & Friction Laws & Harsha \\
%\hline
%8  & Fri 13 Nov 2026 & Laboratory Friction & Carolina \\
%\hline
%9  & Fri 20 Nov 2026 & The Dynamic Rupture Problem & Harsha \\
%\hline
%10 & Fri 04 Dec 2026 & Supershear Rupture & Harsha \\
%\hline
%11 & Fri 11 Dec 2026 & Slow-to-Fast Slip Spectrum & Harsha \\
%\hline
%12 & Fri 18 Dec 2026 & Energy Budget \& Statistical Seismicity & Pierre \\
%\hline
%13 & Fri 08 Jan 2027 & Laboratory Earthquakes & Carolina \& Harsha \\
%\hline
%\end{longtable}
%{\small * Vacances de la Toussaint}
%
%\bigskip

\subsection*{Lecture 1: Continuum Mechanics \& Math Primer}
\begin{itemize}
    \item Why continuum mechanics matters for understanding earthquake sources
    \item Vectors \& tensors --- index notation, coordinate transforms, isotropic tensors
    \item Stress --- traction vector, Cauchy's stress tensor, symmetry, principal stresses
    
    \item Strain --- displacement field, infinitesimal strain tensor, compatibility
    \item Constitutive law --- Hooke's law, Lam\'e parameters, a first look at Navier's equation
    
\end{itemize}

\subsection*{Lecture 2: Elastodynamics I: The Wave Equation}
\begin{itemize}
    \item Quick recap of Navier's equation from last time
    \item Getting to the elastic wave equation from momentum balance and Hooke's law
    \item P \& S waves --- Helmholtz decomposition, wave speeds $\alpha$ and $\beta$
    
    \item Plane waves --- solutions, polarization, particle motion
    \item Energy density \& flux in elastic waves, a first look at geometrical spreading
    
\end{itemize}

\subsection*{Lecture 3: Elastodynamics II: Green's Functions \& Representation Theorem}
\begin{itemize}
    \item Quick recap
    \item The point-force Green's function in a full space --- near, intermediate, and far-field terms
    \item Betti's reciprocal theorem
    
    \item Deriving the representation theorem, and the fault as a displacement discontinuity
    \item The moment tensor and the equivalent body-force picture, the double-couple
    
\end{itemize}

\subsection*{Lecture 4: Far-Field Radiation \& the Seismic Source}
\begin{itemize}
    \item Quick recap of the moment tensor from last time
    \item Far-field P/S radiation patterns from a double couple
    \item Seismic moment $M_0$, and how it relates to magnitude scales
    
    \item The source time function and moment-rate function
    \item The $\omega^2$ model --- corner frequency, estimating static stress drop
    
\end{itemize}

\subsection*{Lecture 5: Kinematic Source Models}
\begin{itemize}
    \item Quick recap
    \item Dislocation theory --- Steketee, Burridge--Knopoff, static dislocation solutions
    \item The Haskell model --- a propagating dislocation source, rise time, rupture velocity, unilateral vs.\ bilateral rupture
    
    \item Finite-fault parameterization --- slip distributions, multiple time-window representations
    \item A first look at inversion --- the forward problem, non-uniqueness, why source inversion is ill-posed
    
\end{itemize}

\subsection*{Lecture 6: Fracture Mechanics Fundamentals}
\begin{itemize}
    \item Quick recap, and why we're moving from kinematics to dynamics
    \item The Griffith energy balance and energy release rate $G$
    \item The crack-tip stress field and stress intensity factor $K$
    
    \item Fracture energy $G_c$ and the process-zone idea
    \item Static vs.\ dynamic fracture criteria, the Griffith criterion
    
\end{itemize}

\subsection*{Lecture 7: Friction Laws}
\begin{itemize}
    \item Quick recap
    \item Coulomb friction --- static/kinetic friction, Byerlee's law
    \item Slip-weakening friction --- Ida (1972), Palmer--Rice, the critical slip distance $D_c$, the breakdown zone
    
    \item Rate-and-state friction --- Dieterich, Ruina, the $(a-b)$ parameter, state evolution laws
    \item Stability --- a linear stability analysis, nucleation size
    
\end{itemize}

\subsection*{Lecture 8: Laboratory Friction}

\subsection*{Lecture 9: The Dynamic Rupture Problem}
\begin{itemize}
    \item Quick recap
    \item Setting up the elastodynamic crack problem and the boundary integral idea
    \item The Kostrov self-similar rupture solution
    
    \item Crack-like vs.\ pulse-like rupture, Heaton's self-healing pulse
    \item The energy budget --- radiated energy, fracture energy, frictional heat, how it all partitions
    
\end{itemize}

\subsection*{Lecture 10: Supershear Rupture}
\begin{itemize}
    \item Quick recap
    \item Sub-Rayleigh vs.\ supershear regimes, the limiting speeds --- Rayleigh, S, P
    \item The Burridge--Andrews mechanism --- daughter-crack nucleation
    
    \item The critical stress ratio $S$ and Mach-cone radiation
    \item What we've observed
    
\end{itemize}

\subsection*{Lecture 11: Slow-to-Fast Slip Spectrum}
\begin{itemize}
    \item Quick recap
    \item What we observe --- slow-slip events, tremor, LFEs, VLFEs
    \item Rate-and-state in the conditionally stable regime
    
    \item Mechanisms --- near-neutral stability, fluid effects, transition zones
    \item Pulling it together: the slip spectrum, from creep to earthquake
    
\end{itemize}

\subsection*{Lecture 12: Energy Budget \& Statistical Seismicity}
\begin{itemize}
    \item Why we're zooming out, from single events to populations of earthquakes
    \item Pulling the energy budget together --- $E_s$, $E_f$, $E_h$, $E_d$, apparent stress, radiation efficiency
    \item How $E_s$ is actually estimated from seismograms --- energy-moment scaling, energy magnitude, and the uncertainties involved
    
    \item The Gutenberg--Richter frequency-magnitude relation and the $b$-value
    \item Omori--Utsu aftershock decay, productivity, a nod to ETAS, and catalog completeness
    
\end{itemize}

\subsection*{Lecture 13: Laboratory Earthquakes}
\begin{itemize}
    \item Photoelasticity experiments
    
    \item Triaxial saw-cut experiments
    
\end{itemize}